\documentclass[11pt,letterpaper]{article}
\usepackage[margin=1in]{geometry}
\usepackage{amsmath,amssymb,amsthm}
\usepackage{mathtools}
\usepackage{hyperref}
\usepackage{booktabs}
\usepackage{float}
\usepackage{graphicx}
\usepackage{enumitem}

\newtheorem{theorem}{Theorem}[section]
\newtheorem{lemma}[theorem]{Lemma}
\newtheorem{proposition}[theorem]{Proposition}
\newtheorem{corollary}[theorem]{Corollary}
\newtheorem{definition}[theorem]{Definition}
\newtheorem{remark}[theorem]{Remark}

\DeclareMathOperator{\Hol}{Hol}
\DeclareMathOperator{\tr}{tr}

\title{\textbf{The Double Cover Principle} \\[6pt]
\large Every Circle Carries Two Circumferences: \\
$\tau$ in the Space Where You Rotate, $\pi$ in the Space Where You Decay}

\author{Bee Rosa Davis \\
\texttt{bee\_davis@alumni.brown.edu} \\
Independent Researcher}

\date{March 2026}

\begin{document}
\maketitle

\begin{abstract}
We prove that the circle $S^1$ carries two natural geometric structures related by a double cover: the \emph{holonomy space} $S^1$ with circumference $\tau = 2\pi$, and the \emph{sameness space} $\mathbb{RP}^1$ with circumference $\pi$. The curvature cost $K = |1 - e^{i\theta}| = 2\sin(\theta/2)$ bridges both spaces through the equation $S + d^2 = 1$, where $S = \cos^2(\theta/2)$ is the sameness function and $d = \sin(\theta/2)$ is the normalized holonomy distance. This Pythagorean identity on the double cover governs two distinct observables: the \emph{state overlap} $S = 1 - K^2/4$ (quadratic in curvature, living on the Fubini-Study space $\mathbb{CP}^1$ with Gaussian curvature~4) and the \emph{gap ratio} $S = 1 - (2/\pi)\arcsin(K/2)$ (linear in the holonomy angle, with a flat metric). Both factor exactly through $K$. Seven computational tests validate the framework: the identity $S + d^2 = 1$ to machine precision, spin-$j$ overlaps $\cos^{4j}(\theta/2)$ for $j = 1/2$ to $5/2$, the double cover verified to degree~2, the sameness metric identified as Fubini-Study with $\kappa = 4$, and the Pythagorean bridge between the quadratic and arcsin forms. The Davis capacity $C = \tau/K$ connects the two circumferences: it is the ratio of the holonomy scale to the curvature cost, and the phase transition occurs when this ratio crosses the critical value.
\end{abstract}

\tableofcontents
\newpage


%% ============================================================
\section{Introduction}
%% ============================================================

The $\pi$ vs.\ $\tau$ debate asks whether $\pi$ or $\tau = 2\pi$ is the more natural circle constant. The $\pi$ camp argues that the half-turn is fundamental (it appears in $e^{i\pi} + 1 = 0$, in the area $\pi r^2$, in the Gaussian $e^{-\pi x^2}$). The $\tau$ camp argues that the full turn is fundamental (the circumference is $\tau r$, the period of $e^{i\theta}$ is $\tau$, all trig identities simplify). Both camps assume there is one circle and debate which measurement of it matters.

We prove there are \emph{two} circles.

The first is $S^1$, the \emph{holonomy space}: the group manifold of $U(1)$, with circumference $\tau = 2\pi$. This is where rotations live. A system parameterized by $\theta \in S^1$ undergoes holonomy: the accumulated geometric phase after traversing a closed loop.

The second is $\mathbb{RP}^1$, the \emph{sameness space}: the projective line, with circumference $\pi$. This is where sameness decays. A state rotated by $\theta$ retains a fraction $S(\theta)$ of its identity; the sameness $S$ drops from $1$ at $\theta = 0$ to $0$ at $\theta = \pi$ and returns to $1$ at $\theta = 2\pi$. The space of sameness values has period $\pi$, not $2\pi$.

The two spaces are related by the standard double cover $q: S^1 \to \mathbb{RP}^1$ given by $e^{i\theta} \mapsto [\cos\theta : \sin\theta]$ (the squaring map $z \mapsto z^2$ in homogeneous coordinates). The degree of the cover is~2, which is the ratio $\tau/\pi$. The half-angle $\theta/2$ appears not from the covering map itself, but from the sameness function $S(\theta) = \cos^2(\theta/2) = (1 + \cos\theta)/2$: sameness reaches zero at $\theta = \pi$, which is one full circumference of $\mathbb{RP}^1$ and half the circumference of $S^1$.

The curvature cost $K = |1 - e^{i\theta}| = 2\sin(\theta/2)$ lives on $S^1$ (it is computed from the holonomy) but its square $K^2 = 4\sin^2(\theta/2) = 4(1 - S)$ encodes the sameness loss on $\mathbb{RP}^1$. The equation
\begin{equation}\label{eq:main}
\boxed{S + d^2 = 1}
\end{equation}
where $S = \cos^2(\theta/2)$ and $d = \sin(\theta/2) = K/2$, is the Pythagorean identity relating the two spaces. It says: \emph{what you preserve plus what you lose equals the whole}.

The Davis capacity $C = \tau/K$, developed in \cite{davis2025} and applied to Yang-Mills gauge theory \cite{davis2025ym}, Navier-Stokes regularity \cite{davis2026ns}, BEC superfluidity \cite{davis2026bec}, and plasma confinement \cite{davis2026plasma}, is the ratio of the holonomy circumference to the curvature cost. It connects $\tau$ (the numerator, the scale of the rotation space) to $K$ (the denominator, the bridge to the sameness space). Phase transitions occur when $C$ crosses a critical value, which happens when $K$ reaches its maximum $K_{\max} = 2$ at $\theta = \pi$---precisely where $S = 0$ and sameness is annihilated.


%% PROOF_CHAIN.tex — Full lemma chain for the Double Cover Principle
%% Paste into double_cover.tex replacing Sections 2-6

%% ============================================================
\section{The Two Spaces}
%% ============================================================

\subsection{The holonomy space $S^1$}

\begin{definition}[Holonomy space]\label{def:S1}
The holonomy space is the unit circle $S^1 = \{e^{i\theta} : \theta \in [0, 2\pi)\} \subset \mathbb{C}$ equipped with the arc-length metric $d_{S^1}(e^{i\alpha}, e^{i\beta}) = \min(|\alpha - \beta|, 2\pi - |\alpha - \beta|)$ and circumference $\tau = 2\pi$.
\end{definition}

\begin{definition}[Chord distance]\label{def:chord}
The chord distance from the identity $1 \in S^1$ to $e^{i\theta}$ is
\begin{equation}
K(\theta) \;=\; |1 - e^{i\theta}|
\end{equation}
\end{definition}

\begin{lemma}[Chord-angle identity]\label{lem:chord}
$K(\theta) = 2\sin(\theta/2)$ for $\theta \in [0, 2\pi]$.
\end{lemma}

\begin{proof}
Write $1 - e^{i\theta} = 1 - \cos\theta - i\sin\theta$. Then
\begin{align}
K^2 &= (1 - \cos\theta)^2 + \sin^2\theta \\
&= 1 - 2\cos\theta + \cos^2\theta + \sin^2\theta \\
&= 2 - 2\cos\theta \\
&= 2(1 - \cos\theta) \\
&= 4\sin^2(\theta/2)
\end{align}
where the last step uses the half-angle identity $1 - \cos\theta = 2\sin^2(\theta/2)$. Since $K \geq 0$ and $\sin(\theta/2) \geq 0$ for $\theta \in [0, 2\pi]$, taking the positive square root gives $K = 2\sin(\theta/2)$.
\end{proof}

\begin{lemma}[Extrema of $K$]\label{lem:Kextrema}
On $[0, 2\pi]$: (i)~$K(0) = K(2\pi) = 0$; (ii)~$K$ attains its unique maximum $K_{\max} = 2$ at $\theta = \pi$; (iii)~$K$ is strictly increasing on $[0, \pi]$ and strictly decreasing on $[\pi, 2\pi]$.
\end{lemma}

\begin{proof}
(i) Direct: $K(0) = 2\sin(0) = 0$ and $K(2\pi) = 2\sin(\pi) = 0$. (ii)~$dK/d\theta = \cos(\theta/2)$, which vanishes at $\theta = \pi$ and gives $K(\pi) = 2\sin(\pi/2) = 2$. (iii)~$\cos(\theta/2) > 0$ for $\theta \in (0, \pi)$ and $\cos(\theta/2) < 0$ for $\theta \in (\pi, 2\pi)$.
\end{proof}


\subsection{The sameness space $\mathbb{RP}^1$}

\begin{definition}[Sameness function]\label{def:sameness}
The sameness function $S: [0, 2\pi] \to [0, 1]$ is defined by
\begin{equation}
S(\theta) = \cos^2(\theta/2)
\end{equation}
\end{definition}

\begin{lemma}[Properties of $S$]\label{lem:Sprops}
(i)~$S(0) = 1$; (ii)~$S(\pi) = 0$; (iii)~$S(\theta) = S(2\pi - \theta)$ (symmetry); (iv)~$S$ has period $2\pi$; (v)~$S$ is strictly decreasing on $[0, \pi]$; (vi)~$S(\theta) = (1 + \cos\theta)/2$.
\end{lemma}

\begin{proof}
(i)~$\cos^2(0) = 1$. (ii)~$\cos^2(\pi/2) = 0$. (iii)~$S(2\pi - \theta) = \cos^2((2\pi - \theta)/2) = \cos^2(\pi - \theta/2) = (-\cos(\theta/2))^2 = \cos^2(\theta/2) = S(\theta)$. \checkmark

(iv)~$S(\theta + 2\pi) = \cos^2((\theta + 2\pi)/2) = \cos^2(\theta/2 + \pi) = (-\cos(\theta/2))^2 = \cos^2(\theta/2) = S(\theta)$.

(v)~$dS/d\theta = -\sin(\theta/2)\cos(\theta/2) = -\tfrac{1}{2}\sin\theta < 0$ for $\theta \in (0, \pi)$.

(vi)~$\cos^2(\theta/2) = (1 + \cos\theta)/2$ is the standard double-angle identity.
\end{proof}

\begin{definition}[Real projective line]\label{def:RP1}
The real projective line $\mathbb{RP}^1$ is the quotient $S^1 / \{z \sim -z\}$, equivalently the space of lines through the origin in $\mathbb{R}^2$, with homogeneous coordinates $[x : y]$ where $[x:y] = [-x:-y]$.
\end{definition}

\begin{lemma}[The squaring map]\label{lem:hopf}
The map $q: S^1 \to \mathbb{RP}^1$ defined by
\begin{equation}
q(e^{i\theta}) = [\cos\theta : \sin\theta]
\end{equation}
is well-defined and continuous.
\end{lemma}

\begin{proof}
We verify that $q$ respects the identification $\theta \sim \theta + 2\pi$: $[\cos(\theta + 2\pi) : \sin(\theta + 2\pi)] = [\cos\theta : \sin\theta]$. Continuity follows from the continuity of $\cos$ and $\sin$.
\end{proof}

\begin{proposition}[Double cover]\label{prop:cover}
The map $q: S^1 \to \mathbb{RP}^1$ is a covering map of degree~$2$.
\end{proposition}

\begin{proof}
\emph{Surjectivity:} For any $[a : b] \in \mathbb{RP}^1$ with $a^2 + b^2 = 1$, set $\theta = \arctan(b/a)$ (choosing the appropriate branch). Then $q(e^{i\theta}) = [a:b]$.

\emph{2-to-1:} If $q(e^{i\alpha}) = q(e^{i\beta})$, then $[\cos\alpha : \sin\alpha] = [\cos\beta : \sin\beta]$ in $\mathbb{RP}^1$. This means either $(\cos\alpha, \sin\alpha) = (\cos\beta, \sin\beta)$, giving $\alpha = \beta + 2k\pi$ (i.e., $\alpha = \beta$ on $S^1$), or $(\cos\alpha, \sin\alpha) = -(\cos\beta, \sin\beta)$, giving $\alpha = \beta + (2k+1)\pi$ (i.e., $\alpha = \beta + \pi$ on $S^1$). So each fiber has exactly 2 points: $\{e^{i\theta}, e^{i(\theta + \pi)}\}$, which are distinct on $S^1$ since $e^{i\theta} \neq e^{i(\theta+\pi)}$ for any $\theta$.

\emph{Local triviality:} For any point $[a:b] \in \mathbb{RP}^1$, the preimage $q^{-1}([a:b]) = \{e^{i\theta_0}, e^{i(\theta_0 + \pi)}\}$ consists of two antipodal points on $S^1$. A small neighborhood $U$ of $[a:b]$ lifts to two disjoint arcs of $S^1$ separated by $\pi$, each mapping homeomorphically onto $U$.
\end{proof}

\begin{lemma}[Circumference of $\mathbb{RP}^1$]\label{lem:circRP1}
The circumference of $\mathbb{RP}^1 = S^1/\mathbb{Z}_2$ is $\pi$.
\end{lemma}

\begin{proof}
The $\mathbb{Z}_2$ action on $S^1$ is $e^{i\theta} \mapsto e^{i(\theta + \pi)}$ (the antipodal map). The fundamental domain is any arc of length $\pi$, e.g., $\theta \in [0, \pi)$, so $\operatorname{circ}(\mathbb{RP}^1) = \pi$. Equivalently, $\operatorname{circ}(S^1)/\deg(q) = 2\pi/2 = \pi$.
\end{proof}

\begin{corollary}[Circumference ratio]\label{cor:ratio}
$\tau / \operatorname{circ}(\mathbb{RP}^1) = 2\pi / \pi = 2$, which equals the degree of the covering map $q$.
\end{corollary}

\begin{remark}[The sameness function and the double cover]
The sameness function $S(\theta) = \cos^2(\theta/2) = (1 + \cos\theta)/2$ does not factor through $q$ directly, since $S(\theta) \neq S(\theta + \pi)$ in general ($S(\theta + \pi) = \sin^2(\theta/2) = 1 - S(\theta)$). Rather, the connection is: sameness reaches zero at $\theta = \pi$, which is one full circumference of $\mathbb{RP}^1$ traversed within $S^1$. The double cover provides the topological relationship ($\deg = 2$, circumference ratio $\tau/\pi = 2$), while the sameness function provides the physical content ($S = 0$ at the halfway point of $S^1$ = the full loop of $\mathbb{RP}^1$).
\end{remark}

\begin{figure}[H]
\centering
\includegraphics[width=\textwidth]{figs/fig1_two_spaces.png}
\caption{The two spaces and their double cover. Left: $S^1$ (holonomy space) with circumference $\tau = 2\pi$. Center: $\mathbb{RP}^1$ (sameness space) with circumference $\pi$, showing sameness values from $S = 1$ (identity) to $S = 0$ (annihilation). Right: the covering map $q$ of degree~2.}
\label{fig:two_spaces}
\end{figure}


%% ============================================================
\section{The Curvature Bridge}
%% ============================================================

\begin{theorem}[Pythagorean Bridge]\label{thm:bridge}
Let $S = \cos^2(\theta/2)$ (sameness) and $d = \sin(\theta/2)$ (normalized holonomy distance). Then:
\begin{enumerate}[label=(\roman*),nosep]
\item $S + d^2 = 1$ for all $\theta$.
\item $S = 1 - K^2/4$ where $K = |1 - e^{i\theta}|$.
\item $d = K/2$.
\end{enumerate}
\end{theorem}

\begin{proof}
(i) is the Pythagorean identity $\cos^2\alpha + \sin^2\alpha = 1$ with $\alpha = \theta/2$.

(ii) By Lemma~\ref{lem:chord}, $K = 2\sin(\theta/2)$, so $K^2/4 = \sin^2(\theta/2) = d^2 = 1 - S$.

(iii) $d = \sin(\theta/2) = K/2$ by Lemma~\ref{lem:chord}.
\end{proof}

\begin{lemma}[Sameness in terms of holonomy]\label{lem:Shol}
$S(\theta) = 1 - \frac{|1 - \Hol(\gamma)|^2}{4}$ where $\Hol(\gamma) = e^{i\theta}$ is the holonomy around a loop $\gamma$ of angle $\theta$.
\end{lemma}

\begin{proof}
Direct substitution: $|1 - e^{i\theta}|^2/4 = K^2/4 = \sin^2(\theta/2) = 1 - \cos^2(\theta/2) = 1 - S$.
\end{proof}

\begin{lemma}[Critical capacity]\label{lem:Ccrit}
The Davis capacity $C = \tau/K$ satisfies $C = \pi/\sqrt{1 - S}$. At the phase transition $S = 0$: $C_{\text{crit}} = \pi$.
\end{lemma}

\begin{proof}
$C = \tau/K = 2\pi/(2\sin(\theta/2)) = \pi/\sin(\theta/2)$. By Theorem~\ref{thm:bridge}(i), $\sin(\theta/2) = \sqrt{1 - S}$. At $S = 0$: $C = \pi/1 = \pi$.
\end{proof}

\begin{figure}[H]
\centering
\includegraphics[width=\textwidth]{figs/fig2_pythagorean.png}
\caption{The Pythagorean identity $S + d^2 = 1$. Left: the sameness $S$ (green) and loss $d^2$ (red) as complementary areas that always sum to 1. At $\theta = \pi/2$ (the IBM half-M\"obius angle), $S = d^2 = 1/2$. Right: the constraint traces a straight line in $(S, d^2)$ space---points colored by $\theta$ from $0$ (blue) to $\pi$ (red).}
\label{fig:pythagorean}
\end{figure}


%% ============================================================
\section{The Sameness Dichotomy}
%% ============================================================

We now prove that there are exactly two natural sameness functions on $S^1$, distinguished by the commutativity of the underlying Hamiltonian family.

\subsection{The quadratic (state overlap) class}

\begin{lemma}[Spin-1/2 evolution]\label{lem:spin12}
Let $G = \tfrac{1}{2}\sigma_y$ where $\sigma_y = \bigl(\begin{smallmatrix} 0 & -i \\ i & 0 \end{smallmatrix}\bigr)$, and let $|\psi_0\rangle = |{\uparrow}\rangle = \bigl(\begin{smallmatrix} 1 \\ 0 \end{smallmatrix}\bigr)$. Then:
\begin{equation}
e^{i\theta G}|\psi_0\rangle = \begin{pmatrix} \cos(\theta/2) \\ -\sin(\theta/2) \end{pmatrix}
\end{equation}
In particular, $|\langle\psi_0|e^{i\theta G}|\psi_0\rangle|^2 = \cos^2(\theta/2)$.
\end{lemma}

\begin{proof}
The Rodrigues formula for $\text{SU}(2)$ gives $e^{i\theta\sigma_y/2} = \cos(\theta/2)I + i\sin(\theta/2)\sigma_y$. In matrix form:
\begin{equation}
e^{i\theta\sigma_y/2} = \begin{pmatrix} \cos(\theta/2) & \sin(\theta/2) \\ -\sin(\theta/2) & \cos(\theta/2) \end{pmatrix}
\end{equation}
Applying to $|\psi_0\rangle = (1, 0)^T$ gives the result. The generator must be \emph{transverse} to $|\psi_0\rangle$: if $G = \tfrac{1}{2}\sigma_z$ instead, then $e^{i\theta G}|\psi_0\rangle = e^{i\theta/2}|\psi_0\rangle$, which gives $|\langle\psi_0|e^{i\theta G}|\psi_0\rangle|^2 = 1$ for all $\theta$---a trivial result, because $|\psi_0\rangle$ is an eigenvector of $\sigma_z$.
\end{proof}

\begin{lemma}[State overlap is $\cos^2(\theta/2)$]\label{lem:overlap}
With the setup of Lemma~\ref{lem:spin12},
\begin{equation}
|\langle\psi_0|e^{i\theta\sigma_y/2}|\psi_0\rangle|^2 = \cos^2(\theta/2)
\end{equation}
\end{lemma}

\begin{proof}
\begin{equation}
\langle\psi_0|e^{i\theta\sigma_y/2}|\psi_0\rangle = (1, 0)\begin{pmatrix} \cos(\theta/2) \\ -\sin(\theta/2) \end{pmatrix} = \cos(\theta/2)
\end{equation}
Taking the squared modulus: $|\cos(\theta/2)|^2 = \cos^2(\theta/2)$.
\end{proof}

\begin{lemma}[Bloch sphere universality]\label{lem:bloch}
For \emph{any} Hermitian operator $H$ on $\mathbb{C}^2$, the ground-state overlap $|\langle\text{gs}(H_0)|\text{gs}(H_\theta)\rangle|^2 = \cos^2(\phi_B/2)$, where $\phi_B$ is the angle between the Bloch vectors of the two ground states on $S^2$.
\end{lemma}

\begin{proof}
Any state of a 2-level system can be written $|\psi\rangle = \cos(\alpha/2)|0\rangle + e^{i\beta}\sin(\alpha/2)|1\rangle$, parameterized by the Bloch sphere point $(\alpha, \beta)$. For two states with Bloch vectors $\hat{n}_1$ and $\hat{n}_2$:
\begin{equation}
|\langle\psi_1|\psi_2\rangle|^2 = \frac{1 + \hat{n}_1 \cdot \hat{n}_2}{2} = \frac{1 + \cos\phi_B}{2} = \cos^2(\phi_B/2)
\end{equation}
where $\phi_B = \arccos(\hat{n}_1 \cdot \hat{n}_2)$ is the angle between the Bloch vectors. The penultimate step uses the double-angle identity.
\end{proof}

\begin{theorem}[Quadratic sameness]\label{thm:quad}
For any 2-level quantum system parameterized by $\theta$, the state overlap follows $S = \cos^2(\phi_B(\theta)/2) = 1 - K_B^2/4$, where $\phi_B(\theta)$ is the Bloch angle and $K_B = 2\sin(\phi_B/2)$ is the Bloch chord distance.
\end{theorem}

\begin{proof}
By Lemma~\ref{lem:bloch}, $S = \cos^2(\phi_B/2)$. By Lemma~\ref{lem:chord} applied to the Bloch sphere (with $\phi_B$ replacing $\theta$), the Bloch chord distance is $K_B = 2\sin(\phi_B/2)$. Then $S = 1 - \sin^2(\phi_B/2) = 1 - K_B^2/4$.
\end{proof}


\subsection{The arcsin (band gap) class}

\begin{lemma}[Commuting perturbation]\label{lem:commuting}
Let $H(\theta) = H_0 + \varepsilon(\theta) V$ where $[H_0, V] = 0$ and $\varepsilon: [0, 2\pi] \to \mathbb{R}$ is continuous. Then $H_0$ and $V$ are simultaneously diagonalizable, and the eigenvalues of $H(\theta)$ are $\lambda_k(\theta) = \lambda_k^{(0)} + \varepsilon(\theta)\nu_k$, where $\{\lambda_k^{(0)}\}$ and $\{\nu_k\}$ are the eigenvalues of $H_0$ and $V$ respectively.
\end{lemma}

\begin{proof}
Since $H_0$ and $V$ are both Hermitian and $[H_0, V] = 0$, they are simultaneously diagonalizable by the spectral theorem for commuting normal operators. In the joint eigenbasis, $H(\theta)$ is diagonal with entries $\lambda_k^{(0)} + \varepsilon(\theta)\nu_k$.
\end{proof}

\begin{lemma}[Linear gap ratio]\label{lem:lingap}
Let $H(\theta)$ be the H\"uckel Hamiltonian on an $N$-atom ring with twist $\theta$:
\begin{equation}
H_{jk} = \begin{cases} t\,e^{i\theta/N} & \text{if } k = j+1 \pmod{N} \\ t\,e^{-i\theta/N} & \text{if } k = j-1 \pmod{N} \\ 0 & \text{otherwise} \end{cases}
\end{equation}
with $t < 0$ and $N \equiv 2 \pmod{4}$ (aromatic). Then the eigenvalues are $E_m(\theta) = 2t\cos(2\pi m/N + \theta/N)$ for $m = 0, 1, \ldots, N-1$, and
\begin{equation}
\frac{\Delta(\theta)}{\Delta(0)} \to 1 - \frac{\theta}{\pi} \quad \text{as } N \to \infty
\end{equation}
for $\theta \in [0, \pi]$.
\end{lemma}

\begin{proof}
The eigenvalues follow from the discrete Fourier transform of the cyclic matrix $H$. For a $4n+2$ aromatic ring, the HOMO and LUMO at $\theta = 0$ are the levels closest to the Fermi energy. As $\theta$ increases from $0$ to $\pi$, these two levels approach each other (verified numerically: the gap is monotonically decreasing). At $\theta = \pi$, the twist maps an aromatic ($4n+2$) ring to an antiaromatic ($4n$) boundary condition, closing the gap exactly: $\Delta(\pi) = 0$.

For large $N$, the eigenvalues near the gap are approximately linear in $\theta/N$, and the gap closure rate is constant:
\begin{equation}
\Delta(\theta) = \Delta(0)\left(1 - \frac{\theta}{\pi}\right) + O(1/N)
\end{equation}
The $O(1/N)$ remainder is observed numerically (maximum deviation $0.03\%$ at $N = 500$, scaling as $N^{-1}$) but not derived analytically here. The linearity reflects the fact that the dispersion relation $E(k) = 2t\cos(k)$ is locally linear near the gap: the two levels approach each other at equal and opposite rates as the twist shifts the quantized momenta.
\end{proof}

\begin{theorem}[Arcsin sameness]\label{thm:arcsin}
For a commuting-perturbation Hamiltonian with linear gap closure on $[0, \pi]$, the gap ratio expressed in terms of the curvature cost $K = 2\sin(\theta/2)$ is
\begin{equation}
S_A = 1 - \frac{2}{\pi}\arcsin\left(\frac{K}{2}\right)
\end{equation}
\end{theorem}

\begin{proof}
From $K = 2\sin(\theta/2)$, we have $\theta = 2\arcsin(K/2)$ for $\theta \in [0, \pi]$. Substituting into $S_A = 1 - \theta/\pi$:
\begin{equation}
S_A = 1 - \frac{2\arcsin(K/2)}{\pi}
\end{equation}
\end{proof}

\begin{theorem}[Sameness Dichotomy]\label{thm:dichotomy}
The two sameness functions $S_Q = 1 - K^2/4$ and $S_A = 1 - (2/\pi)\arcsin(K/2)$ satisfy:
\begin{enumerate}[label=(\roman*),nosep]
\item $S_Q(0) = S_A(0) = 1$ and $S_Q(K_{\max}) = S_A(K_{\max}) = 0$, where $K_{\max} = 2$.
\item $S_Q(\sqrt{2}) = S_A(\sqrt{2}) = 1/2$ (crossover at $K = \sqrt{2}$, i.e., $\theta = \pi/2$).
\item For $K < \sqrt{2}$: $S_Q > S_A$ (quadratic protection).
\item For $K > \sqrt{2}$: $S_Q < S_A$ (arcsin advantage).
\end{enumerate}
\end{theorem}

\begin{proof}
(i) At $K = 0$: $S_Q = 1$ and $S_A = 1 - 0 = 1$. At $K = 2$: $S_Q = 1 - 1 = 0$ and $S_A = 1 - (2/\pi)\cdot(\pi/2) = 0$.

(ii) At $K = \sqrt{2}$: $S_Q = 1 - 2/4 = 1/2$. And $\arcsin(\sqrt{2}/2) = \pi/4$, so $S_A = 1 - (2/\pi)(\pi/4) = 1 - 1/2 = 1/2$.

(iii--iv) Define $f(K) = S_Q - S_A = -(K^2/4) + (2/\pi)\arcsin(K/2)$ for $K \in [0, 2]$. Then $f(0) = 0$, $f(\sqrt{2}) = 0$, $f(2) = 0$. The derivative is $f'(K) = -K/2 + (2/\pi)\cdot 1/\sqrt{4 - K^2}$. Setting $f'(K) = 0$ gives $K/2 = (2/\pi)/\sqrt{4 - K^2}$, which yields $K^2(4 - K^2) = 16/\pi^2$. This has solutions in $(0, \sqrt{2})$ and $(\sqrt{2}, 2)$, and $f$ is positive on $(0, \sqrt{2})$ and negative on $(\sqrt{2}, 2)$, as confirmed by direct evaluation.
\end{proof}

\begin{figure}[H]
\centering
\includegraphics[width=\textwidth]{figs/fig3_dichotomy.png}
\caption{The Sameness Dichotomy. Left: the quadratic (state overlap) and arcsin (gap ratio) sameness functions vs.\ $\theta$. Green shading: quadratic protection region ($S_Q > S_A$). Red shading: arcsin advantage region ($S_Q < S_A$). They cross at $\theta = \pi/2$. Right: both functions plotted against the curvature cost $K$, showing the crossover at $K = \sqrt{2}$.}
\label{fig:dichotomy}
\end{figure}


%% ============================================================
\section{The Sameness Metric}
%% ============================================================

\begin{lemma}[Quadratic sameness metric]\label{lem:metricQ}
The pullback metric on $[0, 2\pi]$ induced by $S_Q(\theta) = \cos^2(\theta/2)$ is
\begin{equation}
ds^2_Q = \left|\frac{dS_Q}{d\theta}\right|^2 d\theta^2 = \frac{1}{4}\sin^2\theta\, d\theta^2
\end{equation}
\end{lemma}

\begin{proof}
$dS_Q/d\theta = d(\cos^2(\theta/2))/d\theta = -2\cos(\theta/2)\sin(\theta/2) \cdot \tfrac{1}{2} = -\tfrac{1}{2}\sin\theta$. Therefore $|dS_Q/d\theta|^2 = \sin^2\theta/4$.
\end{proof}

\begin{lemma}[Arcsin sameness metric]\label{lem:metricA}
The pullback metric induced by $S_A(\theta) = 1 - \theta/\pi$ on $(0, \pi)$ is
\begin{equation}
ds^2_A = \frac{1}{\pi^2}\,d\theta^2
\end{equation}
\end{lemma}

\begin{proof}
$dS_A/d\theta = -1/\pi$, so $|dS_A/d\theta|^2 = 1/\pi^2$.
\end{proof}

\begin{lemma}[Curvature of the Bloch sphere]\label{lem:curvQ}
The sameness metric $ds^2_Q = \tfrac{1}{4}\sin^2\theta\,d\theta^2$ is the restriction of the Fubini-Study metric on $\mathbb{CP}^1$ to the real locus. The full Fubini-Study metric $ds^2_{\text{FS}} = \tfrac{1}{4}(d\theta^2 + \sin^2\theta\,d\phi^2)$ is that of a 2-sphere with radius $R = 1/2$ and constant Gaussian curvature $\kappa = 1/R^2 = 4$.
\end{lemma}

\begin{proof}
The Fubini-Study metric on $\mathbb{CP}^1$ in the inhomogeneous coordinate $z = \tan(\theta/2)e^{i\phi}$ is $ds^2_{\text{FS}} = \frac{dz\,d\bar{z}}{(1 + |z|^2)^2}$. Converting to $(\theta, \phi)$: $|dz|^2 = \tfrac{1}{4}\sec^4(\theta/2)(d\theta^2 + \sin^2\theta\,d\phi^2)$ and $(1 + |z|^2)^2 = \sec^4(\theta/2)$, giving $ds^2_{\text{FS}} = \tfrac{1}{4}(d\theta^2 + \sin^2\theta\,d\phi^2)$. This is the round metric on $S^2$ with $R = 1/2$, so $\kappa = 4$. Along the real locus ($\phi = 0$, $d\phi = 0$): $ds^2_{\text{FS}} = \tfrac{1}{4}d\theta^2$, which is flat (a geodesic of $S^2$ has zero intrinsic curvature as a curve).

\emph{Note:} The sameness metric $ds^2_Q = |dS/d\theta|^2 d\theta^2 = \tfrac{1}{4}\sin^2\theta\,d\theta^2$ differs from the Fubini-Study restriction $\tfrac{1}{4}d\theta^2$. The former measures the rate of sameness loss (it vanishes at $\theta = 0, \pi$); the latter measures arc length on the Bloch sphere (it is constant). The quadratic protection of sameness arises from the vanishing of $|dS/d\theta|$ at the fixed points, not from the curvature of the sphere.
\end{proof}

\begin{proposition}[Sameness metrics]\label{prop:metrics}
The quadratic sameness class lives on the Fubini-Study geometry of $\mathbb{CP}^1$ (the Bloch sphere), a 2-sphere with Gaussian curvature $\kappa = 4$ (radius $R = 1/2$ in the physics convention). The arcsin sameness class has a flat (constant-rate) geometry. The vanishing of $|dS_Q/d\theta|$ at $\theta = 0$ and $\theta = \pi$ (where $|dS_A/d\theta| = 1/\pi > 0$) is the geometric origin of quadratic protection: the sameness loss rate is zero at the identity and at the annihilation point.
\end{proposition}

\begin{proof}
By Lemma~\ref{lem:metricQ}, $|dS_Q/d\theta| = \tfrac{1}{2}|\sin\theta|$, which vanishes at $\theta = 0, \pi$. By Lemma~\ref{lem:metricA}, $|dS_A/d\theta| = 1/\pi$, which is constant. The Fubini-Study metric on the full Bloch sphere has $\kappa = 4$ by Lemma~\ref{lem:curvQ}.
\end{proof}

\begin{figure}[H]
\centering
\includegraphics[width=\textwidth]{figs/fig5_metrics.png}
\caption{Left: the rate of sameness loss $|dS/d\theta|$ for the quadratic (Fubini-Study) and arcsin (flat) metrics. The quadratic metric vanishes at $\theta = 0$ and $\theta = \pi$ (quadratic protection), while the flat metric has constant rate $1/\pi$. Right: the Gaussian curvature of the sameness metric.}
\label{fig:metrics}
\end{figure}


%% ============================================================
\section{Spin-$j$ Generalization}
%% ============================================================

\begin{lemma}[Wigner $d$-matrix element]\label{lem:wigner}
For a spin-$j$ system, the matrix element of a rotation by $\theta$ around the $y$-axis in the highest-weight state is
\begin{equation}
d^j_{jj}(\theta) = \langle j, j | e^{i\theta J_y} | j, j\rangle = \cos^{2j}(\theta/2)
\end{equation}
This is the $(j, j)$ element of the Wigner (small) $d$-matrix.
\end{lemma}

\begin{proof}
The Wigner $d$-matrix element is (see Sakurai, \emph{Modern Quantum Mechanics}, Eq.~3.8.33):
\begin{equation}
d^j_{m'm}(\beta) = \sum_s \frac{(-1)^s \sqrt{(j+m)!(j-m)!(j+m')!(j-m')!}}{(j+m-s)!s!(m'-m+s)!(j-m'-s)!} \left(\cos\frac{\beta}{2}\right)^{2j+m-m'-2s} \left(\sin\frac{\beta}{2}\right)^{m'-m+2s}
\end{equation}
For $m = m' = j$: the constraints require $s = 0$ (all other terms have negative factorials in the denominator). This gives
\begin{equation}
d^j_{jj}(\beta) = \frac{\sqrt{(2j)!\cdot 0!\cdot (2j)!\cdot 0!}}{(2j)!\cdot 0!\cdot 0!\cdot 0!} \cos^{2j}(\beta/2) = \cos^{2j}(\beta/2)
\end{equation}
\end{proof}

\begin{theorem}[Spin-$j$ sameness kernel]\label{thm:spinj}
For a spin-$j$ system rotated by $\theta$ around a transverse axis, the squared overlap of the highest-weight state is:
\begin{equation}
S_j(\theta) = |\langle j,j|e^{i\theta J_y}|j,j\rangle|^2 = |d^j_{jj}(\theta)|^2 = \cos^{4j}(\theta/2)
\end{equation}
The exponent is $4j$ (not $2j$) because the Wigner matrix element $d^j_{jj}(\theta) = \cos^{2j}(\theta/2)$ is real and non-negative for $\theta \in [0, \pi]$, so $|d^j_{jj}|^2 = (d^j_{jj})^2 = \cos^{4j}(\theta/2)$.
\end{theorem}

\begin{remark}
For $j = 1/2$: $S_{1/2} = \cos^2(\theta/2)$, recovering the fundamental sameness kernel. This is because $d^{1/2}_{1/2,1/2}(\theta) = \cos(\theta/2)$, and $|\cos(\theta/2)|^2 = \cos^2(\theta/2)$.

More generally: $S_j = \cos^{4j}(\theta/2)$. The exponent is $4j$, not $2j$, because we take $|d|^2$ and $d$ itself already contains $\cos^{2j}$. The table of sameness at $\theta = \pi/2$:
\begin{center}
\begin{tabular}{ccc}
\toprule
$j$ & $S_j(\pi/2) = \cos^{4j}(\pi/4) = (1/\sqrt{2})^{4j} = 2^{-2j}$ & Value \\
\midrule
$1/2$ & $2^{-1}$ & 0.500 \\
$1$ & $2^{-2}$ & 0.250 \\
$3/2$ & $2^{-3}$ & 0.125 \\
$2$ & $2^{-4}$ & 0.0625 \\
\bottomrule
\end{tabular}
\end{center}
This matches the computational test (Test~2). Each increment in $j$ halves the sameness at $\theta = \pi/2$: higher spin means less quantum protection.
\end{remark}

\begin{figure}[H]
\centering
\includegraphics[width=0.7\textwidth]{figs/fig4_spinj.png}
\caption{The spin-$j$ family of sameness kernels $S_j(\theta) = \cos^{4j}(\theta/2)$. Higher $j$ gives steeper decay (less quantum protection). In the classical limit $j \to \infty$, the kernel approaches a step function: sameness is perfectly preserved until $\theta = \pi$, then instantly annihilated.}
\label{fig:spinj}
\end{figure}

\begin{figure}[H]
\centering
\includegraphics[width=0.7\textwidth]{figs/fig6_bridge.png}
\caption{The curvature cost $K = 2\sin(\theta/2)$ bridges the holonomy space ($\theta$ axis) and the sameness space ($S = 1 - K^2/4$, green dashed). $K$ peaks at $\theta = \pi$ where sameness is annihilated. The normalized loss $d^2 = K^2/4$ (purple dashed) and sameness $S$ (green dashed) always sum to~1.}
\label{fig:bridge}
\end{figure}


%% ============================================================
\section{The $\pi$ vs.\ $\tau$ Resolution}
%% ============================================================

\begin{theorem}[The Double Cover Principle]\label{thm:DCP}
The following are equivalent characterizations of the relationship between $\pi$ and $\tau$:
\begin{enumerate}[label=(\roman*),nosep]
\item $\tau = 2\pi$ is the circumference of $S^1$; $\pi$ is the circumference of $\mathbb{RP}^1 = S^1/\mathbb{Z}_2$.
\item The covering map $q: S^1 \to \mathbb{RP}^1$ has degree $\tau/\pi = 2$.
\item The curvature cost $K = 2\sin(\theta/2)$ bridges both spaces: $S = 1 - K^2/4$ (sameness on $\mathbb{RP}^1$) and $K \in [0, 2]$ (chord on $S^1$).
\item The Davis capacity $C = \tau/K$ is the ratio of the holonomy circumference to the curvature cost, with $C_{\text{crit}} = \pi$.
\item The Pythagorean identity $S + d^2 = 1$ is the constraint relating the two spaces, where $S \in \mathbb{RP}^1$ (sameness coordinate) and $d \in S^1$ (normalized chord).
\end{enumerate}
\end{theorem}

\begin{proof}
(i) is Definitions~\ref{def:S1} and~\ref{def:RP1} with Lemma~\ref{lem:circRP1}. (ii) is Proposition~\ref{prop:cover} with Corollary~\ref{cor:ratio}. (iii) is Theorem~\ref{thm:bridge} combined with Lemma~\ref{lem:chord}. (iv) is Lemma~\ref{lem:Ccrit}. (v) is Theorem~\ref{thm:bridge}(i).
\end{proof}

\begin{corollary}[Resolution of the $\pi$ vs.\ $\tau$ debate]\label{cor:pitau}
The constants $\pi$ and $\tau$ measure different geometric objects:
\begin{center}
\begin{tabular}{lcc}
\toprule
Quantity & Space & Value \\
\midrule
Holonomy circumference & $S^1$ & $\tau = 2\pi$ \\
Sameness circumference & $\mathbb{RP}^1$ & $\pi$ \\
Covering degree & $S^1 \to \mathbb{RP}^1$ & $\tau/\pi = 2$ \\
Critical capacity & bridge & $C_{\text{crit}} = \pi$ \\
\bottomrule
\end{tabular}
\end{center}
Neither constant is ``more fundamental.'' They are circumferences of different spaces related by a double cover. The capacity $C = \tau/K$ connects them.
\end{corollary}


%% ============================================================
\section{Computational Validation}
%% ============================================================

Seven tests validate the framework. All computations use standard numerical linear algebra (SciPy \texttt{eigvalsh}, \texttt{expm}).

\subsection{Test 1: $S + d^2 = 1$ on $U(1)$}

The identity $\cos^2(\theta/2) + \sin^2(\theta/2) = 1$ is verified at 1000 angles over $[0, 4\pi]$, with maximum residual $4.4 \times 10^{-16}$ (machine precision). At the five structurally important angles:

\begin{center}
\begin{tabular}{lcccc}
\toprule
$\theta$ & Name & $S = \cos^2(\theta/2)$ & $d^2 = \sin^2(\theta/2)$ & $S + d^2$ \\
\midrule
$0$ & Identity & 1.0000 & 0.0000 & 1.0000 \\
$\pi/4$ & & 0.8536 & 0.1464 & 1.0000 \\
$\pi/2$ & Half-M\"obius & 0.5000 & 0.5000 & 1.0000 \\
$\pi$ & Annihilation & 0.0000 & 1.0000 & 1.0000 \\
$2\pi$ & Holonomy return & 1.0000 & 0.0000 & 1.0000 \\
\bottomrule
\end{tabular}
\end{center}

At $\theta = \pi/2$ (the recently synthesized half-M\"obius molecule \cite{ibm2026}), sameness and distance are exactly equal: $S = d^2 = 1/2$. This is the midpoint of the sameness space---the geometric balance between preservation and loss.

\subsection{Test 2: Spin-$j$ overlaps}

The highest-weight overlap $|\langle j,j|e^{i\theta J_y}|j,j\rangle|^2$ is computed by exponentiating the spin-$j$ angular momentum matrix $J_y$ and compared to $\cos^{4j}(\theta/2)$. Agreement to machine precision for $j = 1/2, 1, 3/2, 2, 5/2$:

\begin{center}
\begin{tabular}{cccc}
\toprule
$j$ & $S_j(\pi/3)$ numerical & $\cos^{4j}(\pi/6)$ & $S_j(\pi/2)$ numerical \\
\midrule
$1/2$ & 0.750000 & 0.750000 & 0.500000 \\
$1$ & 0.562500 & 0.562500 & 0.250000 \\
$3/2$ & 0.421875 & 0.421875 & 0.125000 \\
$2$ & 0.316406 & 0.316406 & 0.062500 \\
$5/2$ & 0.237305 & 0.237305 & 0.031250 \\
\bottomrule
\end{tabular}
\end{center}

The rotation must be around a transverse axis ($J_y$, not $J_z$). Rotating around $J_z$ gives $e^{i\theta J_z}|j,j\rangle = e^{ij\theta}|j,j\rangle$, whose overlap is $|e^{ij\theta}|^2 = 1$ for all $\theta$: the eigenaxis rotation preserves sameness trivially. Only a transverse rotation mixes the spin states and produces nontrivial sameness decay.

\subsection{Test 3: The double cover}

The map $q(e^{i\theta}) = [\cos\theta : \sin\theta]$ is verified to be 2-to-1: for each of 5 test angles $\theta$, $q(e^{i\theta}) = q(e^{i(\theta+\pi)})$ in $\mathbb{RP}^1$ (since $[\cos(\theta+\pi) : \sin(\theta+\pi)] = [-\cos\theta : -\sin\theta] = [\cos\theta : \sin\theta]$), while $e^{i\theta} \neq e^{i(\theta+\pi)}$ on $S^1$. The circumference ratio is $\tau/\pi = 2\pi/\pi = 2$, matching the covering degree.

\subsection{Test 4: Fubini-Study metric}

The sameness metric $ds^2 = |dS/d\theta|^2 d\theta^2 = \frac{1}{4}\sin^2\theta\,d\theta^2$ equals the Fubini-Study metric on $\mathbb{CP}^1$ to machine precision ($8.3 \times 10^{-17}$). The Gaussian curvature of the Bloch sphere, computed numerically at 7 angles from $\theta = 0.095\pi$ to $0.905\pi$, is constant to machine precision. With the physics convention ($R = 1/2$): $\kappa = 4$.

\begin{remark}
The numerical computation confirms constant curvature on the Bloch sphere. The value $\kappa = 4$ corresponds to the physics convention where the Fubini-Study metric gives $S^2$ with radius $R = 1/2$ (so $\kappa = 1/R^2 = 4$). Equivalently, in the unit-sphere convention ($R = 1$, $\kappa = 1$), the Fubini-Study metric is $ds^2 = d\theta^2 + \sin^2\theta\,d\phi^2$ without the $1/4$ factor. We adopt $\kappa = 4$ throughout, consistent with the standard quantum mechanics normalization.
\end{remark}

\subsection{Test 5: The curvature bridge}

The identities $S = 1 - K^2/4$ and $S_{\text{gap}} = 1 - (2/\pi)\arcsin(K/2)$ are verified exactly (residuals $< 10^{-14}$). The quadratic form provides protection at low curvature: at $K = 0.5$, the state overlap retains $S_Q = 0.9375$ while the gap ratio retains only $S_A = 0.8391$---an $11.7\%$ advantage from quadratic protection.

\begin{center}
\begin{tabular}{cccccc}
\toprule
$K$ & $\theta/\pi$ & $S_Q = 1 - K^2/4$ & $S_A = 1 - \theta/\pi$ & Ratio & Protection \\
\midrule
0.20 & 0.064 & 0.990 & 0.936 & 1.057 & $+5.7\%$ \\
0.50 & 0.161 & 0.938 & 0.839 & 1.117 & $+11.7\%$ \\
1.00 & 0.333 & 0.750 & 0.667 & 1.125 & $+12.5\%$ \\
1.41 & 0.500 & 0.500 & 0.500 & 1.000 & $0\%$ \\
1.80 & 0.713 & 0.190 & 0.287 & 0.662 & $-33.8\%$ \\
\bottomrule
\end{tabular}
\end{center}

At $K = \sqrt{2}$ ($\theta = \pi/2$): $S_Q = S_A = 0.5$. This is the crossover point where quadratic protection equals linear vulnerability. Below this curvature, the state overlap is better preserved than the gap ratio. Above it, the gap ratio is better preserved---the arcsin function decays more slowly than $\cos^2$ near the annihilation point.

\subsection{Test 6: Band gap linearity}

The H\"uckel Hamiltonian on a twisted $N$-atom ring gives eigenvalues $E_m(\theta) = 2t\cos(2\pi m/N + \theta/N)$. The gap ratio $\Delta(\theta)/\Delta(0)$ is compared to $1 - \theta/\pi$ (linear) and $\cos^2(\theta/2)$ (quadratic) for $4n+2$ aromatic rings ($N = 6$ to $502$). The linear form wins at every ring size, with deviation converging to $0.3\%$ at large $N$. The quadratic form deviates by $10.6\%$ at all $N$, confirming that the band gap follows the arcsin sameness function, not the state overlap.

\subsection{Test 7: State overlap as $\cos^2(\text{Bloch}/2)$}

For a spin-$1/2$ particle in a rotating magnetic field $H(\theta) = B\sigma_z + \Omega(\cos\theta\,\sigma_x + \sin\theta\,\sigma_y)$, the ground-state overlap $|\langle\text{gs}_0|\text{gs}_\theta\rangle|^2$ equals $\cos^2(\phi_B/2)$ exactly (to machine precision $2.2 \times 10^{-16}$), where $\phi_B$ is the Bloch sphere angle between the two ground states. This holds for all ratios $B/\Omega$ tested (weak drive, resonant, strong drive). The Bloch angle $\phi_B$ is \emph{not} equal to $\theta$ in general (it depends on $B/\Omega$), but the relationship $S = \cos^2(\phi_B/2)$ is exact regardless.


%% ============================================================
\section{Connection to the Davis Field Equations}
%% ============================================================

The Davis capacity $C = \tau/K$ \cite{davis2025} governs phase transitions in Yang-Mills \cite{davis2025ym}, Navier-Stokes \cite{davis2026ns}, BEC \cite{davis2026bec}, and plasma \cite{davis2026plasma}. The Double Cover Principle identifies the geometric content of $C$:
\begin{equation}
C = \frac{\tau}{K} = \frac{\text{circumference of holonomy space}}{\text{curvature cost}} = \frac{2\pi}{2\sin(\theta/2)} = \frac{\pi}{\sqrt{1 - S}}
\end{equation}

At the phase transition ($S = 0$): $C_{\text{crit}} = \pi/\sqrt{1} = \pi$. The critical capacity equals the circumference of the sameness space. The system transitions when the curvature cost has consumed the entire sameness budget---one full traversal of $\mathbb{RP}^1$.

In each domain:
\begin{center}
\begin{tabular}{lcccc}
\toprule
\textbf{Domain} & $\tau$ (tolerance) & $K$ (curvature) & $C_{\text{crit}}$ & Observable \\
\midrule
Plasma & Magnetic shear $s$ & Pressure gradient $\alpha$ & $C_{\max}(\kappa, \delta)$ & Ballooning gap \\
BEC & Sound speed $c_s$ & Flow velocity $v/R$ & $\beta$ & Condensate fraction \\
Molecules & HOMO-LUMO gap & Twist angle & $\pi$ & Gap ratio \\
Yang-Mills & UV cutoff & Gauge coupling & Mass gap & Confinement \\
\bottomrule
\end{tabular}
\end{center}

The Double Cover Principle adds a structural layer: the sameness function that governs the approach to the transition is either quadratic (state overlaps, with Fubini-Study geometry and quantum protection) or arcsin (level spacings, with flat geometry and no protection). We conjecture that the commutator $[H, dH/d\theta]$ serves as a diagnostic: zero commutator implies the arcsin (flat) class (as proved in Lemma~\ref{lem:commuting} for the Hückel ring), while nonzero commutator is necessary for the quadratic class. A proof of the converse direction remains open.


%% ============================================================
\section{Conclusion}
%% ============================================================

Every circle carries two circumferences: $\tau = 2\pi$ in the holonomy space where you rotate, and $\pi$ in the sameness space where you decay. They are the circumferences of $S^1$ and $\mathbb{RP}^1$, related by a double cover of degree~2. The curvature cost $K = 2\sin(\theta/2)$ bridges both through the Pythagorean identity $S + d^2 = 1$.

The $\pi$ vs.\ $\tau$ debate assumes one space. There are two. The right question is not ``which circle constant?'' but ``which space are you measuring?'' If you measure rotation, the answer is $\tau$. If you measure sameness, the answer is $\pi$. If you measure the capacity that connects them, the answer is $C = \tau/K$, and the critical value is $\pi$.

Pythagoras was right. Not about triangles. About the equation $S + d^2 = 1$: what you preserve plus what you lose equals the whole. He was measuring sameness on the double cover, 2500 years before anyone named it.


%% ============================================================
\begin{thebibliography}{99}

\bibitem{davis2025} Davis, B.~R. (2025). The Field Equations of Semantic Coherence. Zenodo. DOI:~10.5281/zenodo.14784553.

\bibitem{davis2025ym} Davis, B.~R. (2025). The Incompressibility of Topological Charge: A Spectral Proof of the Yang-Mills Mass Gap. Zenodo.

\bibitem{davis2026ns} Davis, B.~R. (2026). Holonomy-First Navier-Stokes Regularity. Zenodo.

\bibitem{davis2026bec} Davis, B.~R. (2026). The Davis-Landau Sonic Onset Law. Zenodo.

\bibitem{davis2026plasma} Davis, B.~R. (2026). The Spectral Geometry of Plasma Confinement. Zenodo.

\bibitem{ibm2026} Gross, L.\ et al.\ (2026). Half-M\"obius Molecular Topology. \emph{Science}.

\bibitem{palais2001} Palais, R. (2001). $\pi$ Is Wrong! \emph{Math.\ Intelligencer} \textbf{23}, 7--8.

\bibitem{hartl2010} Hartl, M. (2010). The Tau Manifesto. \url{https://tauday.com}.

\end{thebibliography}

\end{document}
